\ELhold

\ELsection{Course Outline}{Course Outline}

\begin{ELunit}

\ELfigure{m06-course-outline}{}

\end{ELunit}

\ELhold

\ELsection{Introduction}{Introduction}

\Needspace{11\baselineskip}
\begin{ELunit}

\begin{itemize}
\tightlist
\item
  The fundamental relations of electrostatics and magnetostatics
\end{itemize}

\end{ELunit}

\begin{ELltr}
\begin{longtable}[c]{@{}
  >{\raggedright\arraybackslash}p{(\columnwidth - 4\tabcolsep) * \real{0.2308}}
  >{\centering\arraybackslash}p{(\columnwidth - 4\tabcolsep) * \real{0.3846}}
  >{\centering\arraybackslash}p{(\columnwidth - 4\tabcolsep) * \real{0.3846}}@{}}
\toprule\noalign{}
\ELtablehead \begin{minipage}[b]{\linewidth}\raggedright
Fundamental Relations
\end{minipage} & \begin{minipage}[b]{\linewidth}\centering
Electrostatic Model
\end{minipage} & \begin{minipage}[b]{\linewidth}\centering
Magnetostatic Model
\end{minipage} \\\noalign{\vskip 4pt}
\midrule\noalign{}
\endhead
\bottomrule\noalign{}
\endlastfoot
Governing equations & \ELim{\nabla\times\vect{E}=0}, \ELim{\nabla\cdot\vect{D}=\rho} & \ELim{\nabla\cdot\vect{B}=0}, \ELim{\nabla\times\vect{H}=\vect{J}} \\\noalign{\vskip 4pt}
Constitutive relations (linear and isotropic media) & \ELim{\vect{D}=\epsilon\vect{E}} & \ELim{\vect{H}=\dfrac1\mu\vect{B}} \\\noalign{\vskip 4pt}
\end{longtable}\end{ELltr}


\begin{ELunit}

\begin{itemize}
\tightlist
\item
  In the static (time-invariant) case, the electric field vectors and the magnetic field vectors form separate, independent pairs.
\item
  We shall see that in the time-varying case a time-varying magnetic field produces an electric field, and vice versa
\end{itemize}

\end{ELunit}

\ELhold

\ELsection{Faraday's Law of Electromagnetic Induction}{Faraday’s Law of Electromagnetic Induction}

\begin{ELunit}

\begin{itemize}
\tightlist
\item
  In 1820 \textbf{Hans Christian Ørsted} showed that a wire carrying an electric current deflects a compass needle.

  \begin{itemize}
  \tightlist
  \item
    The motion of the compass needle was caused by the magnetic force of the magnetic field produced by the electric current.
  \end{itemize}
\item
  After this discovery \textbf{Michael Faraday} reasoned that a magnetic field should, in the same way, have an electric effect and produce a current in a wire.
\item
  After about ten years of experiments he made the discovery known as Faraday's law of electromagnetic induction.
\end{itemize}

\end{ELunit}

\begin{ELdefinition}{Faraday's law of electromagnetic induction}

If the magnetic flux through a closed circuit changes with time, an electromotive force is induced in the circuit. This electromotive force drives a current.

\end{ELdefinition}

\begin{ELjoin}

\begin{itemize}
\tightlist
\item
  The polarity of the electromotive force, and hence the direction of the induced current, can be found from Lenz's law.
\end{itemize}

\begin{ELdefinition}{Lenz's law}

The current induced in a circuit always flows in the direction that opposes the change of the magnetic flux producing it.

\end{ELdefinition}

\end{ELjoin}

\begin{ELjoin}

\begin{itemize}
\tightlist
\item
  The mathematical statement of Faraday's law of electromagnetic induction
\end{itemize}

\begin{ELimportant}{}

\ELdisp{v_{emf}\ELbr =-N\frac{d\Phi}{dt}\ELbr =-N\frac{d}{dt}\int_S\vect{B}\cdot d\vect{s}}

\end{ELimportant}

\end{ELjoin}

\begin{ELunit}

\begin{itemize}
\tightlist
\item
  The definition of the electromotive force induced in a circuit with the closed contour \ELim{C}
\end{itemize}

\ELdisp{v_{emf}\triangleq\oint_C\vect{E}\cdot\dif\vect{\ell}}

\begin{itemize}
\tightlist
\item
  For a stationary circuit
\end{itemize}

\ELdisp{v_{emf}\ELbr =\oint_C\vect{E}\cdot\dif\vect{\ell}\ELbr =-\int_S\frac{\partial\vect{B}}{\partial t}\cdot d\vect{s}\quad\ELbr \ELbr \Longrightarrow\quad\ELbr \int_S(\nabla\times\vect{E})\cdot d\vect{s}\ELbr =-\int_S\frac{\partial\vect{B}}{\partial t}\cdot d\vect{s}}

\end{ELunit}

\begin{ELjoin}

\begin{itemize}
\tightlist
\item
  The last step uses Stokes's theorem.
\end{itemize}

\begin{ELimportant}{}

\ELdisp{\nabla\times\vect{E}\ELbr =-\frac{\partial\vect{B}}{\partial t}}

\end{ELimportant}

\end{ELjoin}

\begin{ELunit}

\begin{itemize}
\tightlist
\item
  In the time-varying case the electric field is not conservative \ELim{\ELto} \ELim{\nabla\times\vect{E}\neq0}
\item
  In the time-varying case the scalar potential cannot be defined as \ELim{V=-\int\vect{E}\cdot\dif\vect{\ell}}, because the potential difference between two points would depend on the path of integration and would not have a fixed value.
\item
  At low frequencies \ELim{\ELto} \ELim{\nabla\times\vect{E}\simeq0} \ELim{\ELto} the scalar potential can be defined (lumped-circuit theory)
\end{itemize}

\end{ELunit}

\ELnextnum{6-1}

\begin{ELexample}{}

A circular loop of \ELim{N} turns of conducting wire lies in the \ELim{xy}-plane with its centre at the origin, in a magnetic field given by

\ELdisp{\vect{B}\ELbr =\uvec{z}B_0\cos\left(\frac{\pi r}{2b}\right)\sin(\omega t)}

Here \ELim{b} is the radius of the loop and \ELim{\omega} the angular frequency. Find the induced electromotive force.

\begin{ELsolution}

\ELdisp{\begin{aligned}\Phi=\int_S\vect{B}\cdot d\vect{s}&=\int_0^{2\pi}\!\!\int_0^b\left[B_0\cos\left(\frac{\pi r}{2b}\right)\sin(\omega t)\,\uvec{z}\right]\cdot\left(r\,dr\,d\phi\,\uvec{z}\right)\\&=\frac{8b^2}{\pi}\left(\frac\pi2-1\right)B_0\sin(\omega t)\end{aligned}}
\ELdisp{v\ELbr =-N\frac{d\Phi}{dt}\ELbr =-\frac{8Nb^2}{\pi}\left(\frac\pi2-1\right)B_0\omega\cos(\omega t)}

\end{ELsolution}

\end{ELexample}

\ELhold

\ELsection{Displacement Current}{Displacement Current}

\begin{ELunit}

\begin{itemize}
\tightlist
\item
  Ampère's circuital law for the static magnetic field
\end{itemize}

\ELdisp{\nabla\times\vect{H}\ELbr =\vect{J}}

\begin{itemize}
\tightlist
\item
  \textbf{James Clerk Maxwell} doubted this equation, because
\end{itemize}

\ELdisp{\nabla\cdot(\nabla\times\vect{H})\ELbr =0\ELbr =\nabla\cdot\vect{J}}

\begin{itemize}
\tightlist
\item
  which is not consistent with the equation of continuity
\end{itemize}

\ELdisp{\nabla\cdot\vect{J}\ELbr =-\frac{\partial\rho}{\partial t}}

\begin{itemize}
\tightlist
\item
  To remove this inconsistency Maxwell modified the relation as follows
\end{itemize}

\ELdisp{\nabla\cdot(\nabla\times\vect{H})\ELbr =0\ELbr =\nabla\cdot\vect{J}+\frac{\partial\rho}{\partial t}\quad\ELbr \xrightarrow{\ \nabla\cdot\vect{D}=\rho\ }\quad\ELbr \nabla\cdot(\nabla\times\vect{H})\ELbr =\nabla\cdot\left(\vect{J}+\frac{\partial\vect{D}}{\partial t}\right)}

\end{ELunit}

\begin{ELimportant}{}

\ELdisp{\nabla\times\vect{H}\ELbr =\vect{J}+\frac{\partial\vect{D}}{\partial t}}

\end{ELimportant}

\begin{ELunit}

\begin{itemize}
\tightlist
\item
  It can easily be shown that \ELim{\partial\vect{D}/\partial t} has the dimension of a current density. It is called the \textbf{displacement current density}.
\item
  The modified form of Ampère's circuital law states that a time-varying electric field produces a magnetic field even when there is no current.
\end{itemize}

\end{ELunit}

\ELhold

\ELsection{Maxwell's Equations}{Maxwell’s Equations}

\Needspace{15\baselineskip}
\begin{ELunit}

\begin{itemize}
\tightlist
\item
  The equations governing time-varying electric and magnetic fields, or \textbf{Maxwell's equations}
\end{itemize}

\end{ELunit}

\begin{ELltr}
\begin{longtable}[c]{@{}
  >{\raggedright\arraybackslash}p{(\columnwidth - 4\tabcolsep) * \real{0.3333}}
  >{\raggedright\arraybackslash}p{(\columnwidth - 4\tabcolsep) * \real{0.3333}}
  >{\raggedright\arraybackslash}p{(\columnwidth - 4\tabcolsep) * \real{0.3333}}@{}}
\toprule\noalign{}
\ELtablehead \begin{minipage}[b]{\linewidth}\raggedright
Differential Form
\end{minipage} & \begin{minipage}[b]{\linewidth}\raggedright
Integral Form
\end{minipage} & \begin{minipage}[b]{\linewidth}\raggedright
Significance
\end{minipage} \\\noalign{\vskip 4pt}
\midrule\noalign{}
\endhead
\bottomrule\noalign{}
\endlastfoot
\ELim{\nabla\times\vect{E}=-\dfrac{\partial\vect{B}}{\partial t}} & \ELim{\oint_C\vect{E}\cdot\dif\vect{\ell}=-\dfrac{d\Phi}{dt}} & Faraday's law \\\noalign{\vskip 4pt}
\ELim{\nabla\times\vect{H}=\vect{J}+\dfrac{\partial\vect{D}}{\partial t}} & \ELim{\oint_C\vect{H}\cdot\dif\vect{\ell}=I+\int_S\dfrac{\partial\vect{D}}{\partial t}\cdot d\vect{s}} & Ampère's circuital law \\\noalign{\vskip 4pt}
\ELim{\nabla\cdot\vect{D}=\rho} & \ELim{\oint_S\vect{D}\cdot d\vect{s}=Q} & Gauss's law \\\noalign{\vskip 4pt}
\ELim{\nabla\cdot\vect{B}=0} & \ELim{\oint_S\vect{B}\cdot d\vect{s}=0} & No isolated magnetic charge \\\noalign{\vskip 4pt}
\end{longtable}\end{ELltr}

